% =====================================================================
% SI table | Prior-referenced comparison of TEXAS and BAYSPAR
%
% Numbers generated by notebooks/manuscripts/SI03_paleo_showcases_modelswitch.ipynb
% (final cell, "SI table | prior-referenced comparison"), t0shift / GHEB arm.
%
% Preamble requirements:
%   \usepackage{amsmath,tabularx,booktabs,multirow,gensymb,textcomp}
%
% Layout: panels A and B share the SAME fixed widths for the first three
% columns (Record / Interval / n), set by \colRecord, \colInterval and \colN
% below, so the two panels line up column-for-column. All slack is absorbed
% AFTER column 3 by \extracolsep{\fill}, which is why both panels fill
% \linewidth without disturbing the shared columns. Change a width in one
% place and both panels follow.
%
% Font: the table body is set in TeX Gyre Heros Condensed (qhvc). Every
% NUMBER is therefore typeset in text mode, not math mode -- math mode would
% fall back to the serif math font and break the sans setting. Symbols that
% are genuinely variables (n, mu, sigma, p, z, Phi) stay in math mode, where
% italic serif is correct. To make those sans as well, add to the preamble:
%     \usepackage[italic]{mathastext}   % after the sans font is made default
%
% Use with:  % =====================================================================
% SI table | Prior-referenced comparison of TEXAS and BAYSPAR
%
% Numbers generated by notebooks/manuscripts/SI03_paleo_showcases_modelswitch.ipynb
% (final cell, "SI table | prior-referenced comparison"), t0shift / GHEB arm.
%
% Preamble requirements:
%   \usepackage{amsmath,tabularx,booktabs,multirow,gensymb,textcomp}
%
% Layout: panels A and B share the SAME fixed widths for the first three
% columns (Record / Interval / n), set by \colRecord, \colInterval and \colN
% below, so the two panels line up column-for-column. All slack is absorbed
% AFTER column 3 by \extracolsep{\fill}, which is why both panels fill
% \linewidth without disturbing the shared columns. Change a width in one
% place and both panels follow.
%
% Font: the table body is set in TeX Gyre Heros Condensed (qhvc). Every
% NUMBER is therefore typeset in text mode, not math mode -- math mode would
% fall back to the serif math font and break the sans setting. Symbols that
% are genuinely variables (n, mu, sigma, p, z, Phi) stay in math mode, where
% italic serif is correct. To make those sans as well, add to the preamble:
%     \usepackage[italic]{mathastext}   % after the sans font is made default
%
% Use with:  % =====================================================================
% SI table | Prior-referenced comparison of TEXAS and BAYSPAR
%
% Numbers generated by notebooks/manuscripts/SI03_paleo_showcases_modelswitch.ipynb
% (final cell, "SI table | prior-referenced comparison"), t0shift / GHEB arm.
%
% Preamble requirements:
%   \usepackage{amsmath,tabularx,booktabs,multirow,gensymb,textcomp}
%
% Layout: panels A and B share the SAME fixed widths for the first three
% columns (Record / Interval / n), set by \colRecord, \colInterval and \colN
% below, so the two panels line up column-for-column. All slack is absorbed
% AFTER column 3 by \extracolsep{\fill}, which is why both panels fill
% \linewidth without disturbing the shared columns. Change a width in one
% place and both panels follow.
%
% Font: the table body is set in TeX Gyre Heros Condensed (qhvc). Every
% NUMBER is therefore typeset in text mode, not math mode -- math mode would
% fall back to the serif math font and break the sans setting. Symbols that
% are genuinely variables (n, mu, sigma, p, z, Phi) stay in math mode, where
% italic serif is correct. To make those sans as well, add to the preamble:
%     \usepackage[italic]{mathastext}   % after the sans font is made default
%
% Use with:  % =====================================================================
% SI table | Prior-referenced comparison of TEXAS and BAYSPAR
%
% Numbers generated by notebooks/manuscripts/SI03_paleo_showcases_modelswitch.ipynb
% (final cell, "SI table | prior-referenced comparison"), t0shift / GHEB arm.
%
% Preamble requirements:
%   \usepackage{amsmath,tabularx,booktabs,multirow,gensymb,textcomp}
%
% Layout: panels A and B share the SAME fixed widths for the first three
% columns (Record / Interval / n), set by \colRecord, \colInterval and \colN
% below, so the two panels line up column-for-column. All slack is absorbed
% AFTER column 3 by \extracolsep{\fill}, which is why both panels fill
% \linewidth without disturbing the shared columns. Change a width in one
% place and both panels follow.
%
% Font: the table body is set in TeX Gyre Heros Condensed (qhvc). Every
% NUMBER is therefore typeset in text mode, not math mode -- math mode would
% fall back to the serif math font and break the sans setting. Symbols that
% are genuinely variables (n, mu, sigma, p, z, Phi) stay in math mode, where
% italic serif is correct. To make those sans as well, add to the preamble:
%     \usepackage[italic]{mathastext}   % after the sans font is made default
%
% Use with:  \input{SI_table_prior_referenced}
% =====================================================================

% Fallbacks so the table compiles even without textcomp.
\providecommand{\textminus}{-}
% Nitrate scenario label, in text mode so it follows the table font.
\providecommand{\nitrate}{NO\textsubscript{3}\textsuperscript{\textminus}}
% Shared geometry of the first three columns -- panels A and B both use these.
\providecommand{\colRecord}{2.55cm}
\providecommand{\colInterval}{2.20cm}
\providecommand{\colN}{0.45cm}
% NB: the column tokens themselves must be written out in each panel --
% \array does not expand a macro in the preamble-token position. Only the
% p{} widths above are shared, and they are the only thing to edit.

\begin{table}[h]\centering\footnotesize
\fontfamily{qhvc}\selectfont
\caption{Prior-referenced comparison of TEXAS and BAYSPAR.}

% ---- (A) temperature: did the data move the posterior off the prior? --------
\textbf{(A) Posterior median temperature (\degree C), with the shift from the prior in units of $\sigma_{\text{prior}}$}\\[2pt]

\begin{tabular*}{\linewidth}{@{}>{\raggedright\arraybackslash}p{\colRecord} >{\raggedright\arraybackslash}p{\colInterval} >{\centering\arraybackslash}p{\colN} @{\extracolsep{\fill}} c c c c@{}}
\toprule
\textbf{Record} & \textbf{Interval} & $n$ & $\mu_{\text{T,prior}}$ & \textbf{TEXAS} & \textbf{TEXAS} & \textbf{BAYSPAR} \\
 & & & & (\nitrate=10) & (\nitrate=0.1) & \\
\midrule
\multirow{3}{=}{South Dover Bridge}
 & post-PETM & 3  & 30 & 30.2 (+0.02) & 30.2 (+0.02) & 32.5 (+0.25) \\
 & PETM body & 18 & 35 & 35.5 (+0.05) & 32.9 (\textminus0.21) & 39.7 (+0.47) \\
 & pre-PETM  & 14 & 30 & 28.0 (\textminus0.20) & 28.1 (\textminus0.19) & 27.7 (\textminus0.23) \\
\midrule
\multirow{3}{=}{ODP~959}
 & post-PETM & 4  & 30 & 36.5 (+0.65) & 36.4 (+0.64) & 38.1 (+0.81) \\
 & PETM body & 20 & 35 & 42.3 (+0.73) & 40.3 (+0.53) & 43.1 (+0.81) \\
 & pre-PETM  & 16 & 30 & 34.8 (+0.48) & 34.8 (+0.48) & 35.9 (+0.59) \\
\midrule
ODP~1259 & upper asymptote & 59 & 40 & 40.6 (+0.06) & --- & 41.1 (+0.11) \\
Co1010   & lower asymptote & 36 &  5 &  3.0 (\textminus0.20) & --- &  5.8 (+0.08) \\
\bottomrule
\end{tabular*}

\vspace{8pt}
% ---- (B) width -------------------------------------------------------------
\textbf{(B) One-sided posterior width (\degree C) and prior variance removed}\\[2pt]

\begin{tabular*}{\linewidth}{@{}>{\raggedright\arraybackslash}p{\colRecord} >{\raggedright\arraybackslash}p{\colInterval} >{\centering\arraybackslash}p{\colN} @{\extracolsep{\fill}} l c c c c@{}}
\toprule
\textbf{Record} & \textbf{Interval} & $n$ & \textbf{Method}
 & $\sigma_-$ & $\sigma_+$ & $\bar{\sigma}$ & \textbf{resolved} \\
\midrule
\multirow{9}{=}{South Dover Bridge}
 & \multirow{3}{*}{post-PETM} & \multirow{3}{*}{3}
   & TEXAS (\nitrate=10)  & 1.95 & 1.93 & 1.94 & 96\% \\
 & & & TEXAS (\nitrate=0.1) & 1.94 & 1.89 & 1.91 & 96\% \\
 & & & BAYSPAR              & 3.95 & 4.60 & 4.28 & 82\% \\
\cmidrule(l){2-8}
 & \multirow{3}{*}{PETM body} & \multirow{3}{*}{18}
   & TEXAS (\nitrate=10)  & 2.17 & 2.80 & 2.49 & 94\% \\
 & & & TEXAS (\nitrate=0.1) & 2.20 & 3.05 & 2.62 & 93\% \\
 & & & BAYSPAR              & 4.61 & 5.51 & 5.06 & 74\% \\
\cmidrule(l){2-8}
 & \multirow{3}{*}{pre-PETM} & \multirow{3}{*}{14}
   & TEXAS (\nitrate=10)  & 2.09 & 1.93 & 2.01 & 96\% \\
 & & & TEXAS (\nitrate=0.1) & 2.03 & 1.95 & 1.99 & 96\% \\
 & & & BAYSPAR              & 3.78 & 3.98 & 3.88 & 85\% \\
\midrule
\multirow{9}{=}{ODP~959}
 & \multirow{3}{*}{post-PETM} & \multirow{3}{*}{4}
   & TEXAS (\nitrate=10)  & 2.02 & 2.49 & 2.25 & 95\% \\
 & & & TEXAS (\nitrate=0.1) & 2.07 & 2.57 & 2.32 & 95\% \\
 & & & BAYSPAR              & 4.11 & 4.24 & 4.18 & 83\% \\
\cmidrule(l){2-8}
 & \multirow{3}{*}{PETM body} & \multirow{3}{*}{20}
   & TEXAS (\nitrate=10)  & 3.71 & 5.99 & 4.85 & 76\% \\
 & & & TEXAS (\nitrate=0.1) & 4.09 & 6.73 & 5.41 & 71\% \\
 & & & BAYSPAR              & 4.70 & 4.98 & 4.84 & 77\% \\
\cmidrule(l){2-8}
 & \multirow{3}{*}{pre-PETM} & \multirow{3}{*}{16}
   & TEXAS (\nitrate=10)  & 1.94 & 2.22 & 2.08 & 96\% \\
 & & & TEXAS (\nitrate=0.1) & 1.96 & 2.21 & 2.09 & 96\% \\
 & & & BAYSPAR              & 4.01 & 4.12 & 4.07 & 83\% \\
\midrule
\multirow{2}{=}{ODP~1259} & \multirow{2}{*}{upper asymptote} & \multirow{2}{*}{59}
   & TEXAS (\nitrate=10)  & 3.75 & 7.00 & 5.37 & 71\% \\
 & & & BAYSPAR              & 4.98 & 5.14 & 5.06 & 74\% \\
\midrule
\multirow{2}{=}{Co1010} & \multirow{2}{*}{lower asymptote} & \multirow{2}{*}{36}
   & TEXAS (\nitrate=10)  & 7.33 & 5.85 & 6.59 & 57\% \\
 & & & BAYSPAR              & 5.05 & 4.67 & 4.86 & 76\% \\
\bottomrule
\end{tabular*}

\vspace{2pt}
\begin{minipage}{0.95\linewidth}\footnotesize
Note: $\sigma_-$ and $\sigma_+$ are the one-sided posterior widths below and above the median, $(p_{50}-p_{16})/z$ and $(p_{84}-p_{50})/z$ with $z=\Phi^{-1}(0.84)$; $\bar{\sigma}$ is their mean, the symmetric-equivalent width. \textbf{\emph{Resolved}} is $1-\bar{\sigma}^{2}/\sigma_{\text{prior}}^{2}$, the fraction of the prior variance eliminated by the GDGT data: 100\% means the proxy, not the prior, sets the estimate; 0\% means the posterior is as wide as the prior and the data added nothing. All methods use the same per-sample prior mean $\mu_{\text{T,prior}}$ and $\sigma_{\text{prior}}=10$\,\degree C, and within each interval are averaged over the same samples. \nitrate=0.1 is a sensitivity scenario applied only inside the PETM interval, so it is not run for the two asymptote records.
\end{minipage}
\end{table}

% =====================================================================

% Fallbacks so the table compiles even without textcomp.
\providecommand{\textminus}{-}
% Nitrate scenario label, in text mode so it follows the table font.
\providecommand{\nitrate}{NO\textsubscript{3}\textsuperscript{\textminus}}
% Shared geometry of the first three columns -- panels A and B both use these.
\providecommand{\colRecord}{2.55cm}
\providecommand{\colInterval}{2.20cm}
\providecommand{\colN}{0.45cm}
% NB: the column tokens themselves must be written out in each panel --
% \array does not expand a macro in the preamble-token position. Only the
% p{} widths above are shared, and they are the only thing to edit.

\begin{table}[h]\centering\footnotesize
\fontfamily{qhvc}\selectfont
\caption{Prior-referenced comparison of TEXAS and BAYSPAR.}

% ---- (A) temperature: did the data move the posterior off the prior? --------
\textbf{(A) Posterior median temperature (\degree C), with the shift from the prior in units of $\sigma_{\text{prior}}$}\\[2pt]

\begin{tabular*}{\linewidth}{@{}>{\raggedright\arraybackslash}p{\colRecord} >{\raggedright\arraybackslash}p{\colInterval} >{\centering\arraybackslash}p{\colN} @{\extracolsep{\fill}} c c c c@{}}
\toprule
\textbf{Record} & \textbf{Interval} & $n$ & $\mu_{\text{T,prior}}$ & \textbf{TEXAS} & \textbf{TEXAS} & \textbf{BAYSPAR} \\
 & & & & (\nitrate=10) & (\nitrate=0.1) & \\
\midrule
\multirow{3}{=}{South Dover Bridge}
 & post-PETM & 3  & 30 & 30.2 (+0.02) & 30.2 (+0.02) & 32.5 (+0.25) \\
 & PETM body & 18 & 35 & 35.5 (+0.05) & 32.9 (\textminus0.21) & 39.7 (+0.47) \\
 & pre-PETM  & 14 & 30 & 28.0 (\textminus0.20) & 28.1 (\textminus0.19) & 27.7 (\textminus0.23) \\
\midrule
\multirow{3}{=}{ODP~959}
 & post-PETM & 4  & 30 & 36.5 (+0.65) & 36.4 (+0.64) & 38.1 (+0.81) \\
 & PETM body & 20 & 35 & 42.3 (+0.73) & 40.3 (+0.53) & 43.1 (+0.81) \\
 & pre-PETM  & 16 & 30 & 34.8 (+0.48) & 34.8 (+0.48) & 35.9 (+0.59) \\
\midrule
ODP~1259 & upper asymptote & 59 & 40 & 40.6 (+0.06) & --- & 41.1 (+0.11) \\
Co1010   & lower asymptote & 36 &  5 &  3.0 (\textminus0.20) & --- &  5.8 (+0.08) \\
\bottomrule
\end{tabular*}

\vspace{8pt}
% ---- (B) width -------------------------------------------------------------
\textbf{(B) One-sided posterior width (\degree C) and prior variance removed}\\[2pt]

\begin{tabular*}{\linewidth}{@{}>{\raggedright\arraybackslash}p{\colRecord} >{\raggedright\arraybackslash}p{\colInterval} >{\centering\arraybackslash}p{\colN} @{\extracolsep{\fill}} l c c c c@{}}
\toprule
\textbf{Record} & \textbf{Interval} & $n$ & \textbf{Method}
 & $\sigma_-$ & $\sigma_+$ & $\bar{\sigma}$ & \textbf{resolved} \\
\midrule
\multirow{9}{=}{South Dover Bridge}
 & \multirow{3}{*}{post-PETM} & \multirow{3}{*}{3}
   & TEXAS (\nitrate=10)  & 1.95 & 1.93 & 1.94 & 96\% \\
 & & & TEXAS (\nitrate=0.1) & 1.94 & 1.89 & 1.91 & 96\% \\
 & & & BAYSPAR              & 3.95 & 4.60 & 4.28 & 82\% \\
\cmidrule(l){2-8}
 & \multirow{3}{*}{PETM body} & \multirow{3}{*}{18}
   & TEXAS (\nitrate=10)  & 2.17 & 2.80 & 2.49 & 94\% \\
 & & & TEXAS (\nitrate=0.1) & 2.20 & 3.05 & 2.62 & 93\% \\
 & & & BAYSPAR              & 4.61 & 5.51 & 5.06 & 74\% \\
\cmidrule(l){2-8}
 & \multirow{3}{*}{pre-PETM} & \multirow{3}{*}{14}
   & TEXAS (\nitrate=10)  & 2.09 & 1.93 & 2.01 & 96\% \\
 & & & TEXAS (\nitrate=0.1) & 2.03 & 1.95 & 1.99 & 96\% \\
 & & & BAYSPAR              & 3.78 & 3.98 & 3.88 & 85\% \\
\midrule
\multirow{9}{=}{ODP~959}
 & \multirow{3}{*}{post-PETM} & \multirow{3}{*}{4}
   & TEXAS (\nitrate=10)  & 2.02 & 2.49 & 2.25 & 95\% \\
 & & & TEXAS (\nitrate=0.1) & 2.07 & 2.57 & 2.32 & 95\% \\
 & & & BAYSPAR              & 4.11 & 4.24 & 4.18 & 83\% \\
\cmidrule(l){2-8}
 & \multirow{3}{*}{PETM body} & \multirow{3}{*}{20}
   & TEXAS (\nitrate=10)  & 3.71 & 5.99 & 4.85 & 76\% \\
 & & & TEXAS (\nitrate=0.1) & 4.09 & 6.73 & 5.41 & 71\% \\
 & & & BAYSPAR              & 4.70 & 4.98 & 4.84 & 77\% \\
\cmidrule(l){2-8}
 & \multirow{3}{*}{pre-PETM} & \multirow{3}{*}{16}
   & TEXAS (\nitrate=10)  & 1.94 & 2.22 & 2.08 & 96\% \\
 & & & TEXAS (\nitrate=0.1) & 1.96 & 2.21 & 2.09 & 96\% \\
 & & & BAYSPAR              & 4.01 & 4.12 & 4.07 & 83\% \\
\midrule
\multirow{2}{=}{ODP~1259} & \multirow{2}{*}{upper asymptote} & \multirow{2}{*}{59}
   & TEXAS (\nitrate=10)  & 3.75 & 7.00 & 5.37 & 71\% \\
 & & & BAYSPAR              & 4.98 & 5.14 & 5.06 & 74\% \\
\midrule
\multirow{2}{=}{Co1010} & \multirow{2}{*}{lower asymptote} & \multirow{2}{*}{36}
   & TEXAS (\nitrate=10)  & 7.33 & 5.85 & 6.59 & 57\% \\
 & & & BAYSPAR              & 5.05 & 4.67 & 4.86 & 76\% \\
\bottomrule
\end{tabular*}

\vspace{2pt}
\begin{minipage}{0.95\linewidth}\footnotesize
Note: $\sigma_-$ and $\sigma_+$ are the one-sided posterior widths below and above the median, $(p_{50}-p_{16})/z$ and $(p_{84}-p_{50})/z$ with $z=\Phi^{-1}(0.84)$; $\bar{\sigma}$ is their mean, the symmetric-equivalent width. \textbf{\emph{Resolved}} is $1-\bar{\sigma}^{2}/\sigma_{\text{prior}}^{2}$, the fraction of the prior variance eliminated by the GDGT data: 100\% means the proxy, not the prior, sets the estimate; 0\% means the posterior is as wide as the prior and the data added nothing. All methods use the same per-sample prior mean $\mu_{\text{T,prior}}$ and $\sigma_{\text{prior}}=10$\,\degree C, and within each interval are averaged over the same samples. \nitrate=0.1 is a sensitivity scenario applied only inside the PETM interval, so it is not run for the two asymptote records.
\end{minipage}
\end{table}

% =====================================================================

% Fallbacks so the table compiles even without textcomp.
\providecommand{\textminus}{-}
% Nitrate scenario label, in text mode so it follows the table font.
\providecommand{\nitrate}{NO\textsubscript{3}\textsuperscript{\textminus}}
% Shared geometry of the first three columns -- panels A and B both use these.
\providecommand{\colRecord}{2.55cm}
\providecommand{\colInterval}{2.20cm}
\providecommand{\colN}{0.45cm}
% NB: the column tokens themselves must be written out in each panel --
% \array does not expand a macro in the preamble-token position. Only the
% p{} widths above are shared, and they are the only thing to edit.

\begin{table}[h]\centering\footnotesize
\fontfamily{qhvc}\selectfont
\caption{Prior-referenced comparison of TEXAS and BAYSPAR.}

% ---- (A) temperature: did the data move the posterior off the prior? --------
\textbf{(A) Posterior median temperature (\degree C), with the shift from the prior in units of $\sigma_{\text{prior}}$}\\[2pt]

\begin{tabular*}{\linewidth}{@{}>{\raggedright\arraybackslash}p{\colRecord} >{\raggedright\arraybackslash}p{\colInterval} >{\centering\arraybackslash}p{\colN} @{\extracolsep{\fill}} c c c c@{}}
\toprule
\textbf{Record} & \textbf{Interval} & $n$ & $\mu_{\text{T,prior}}$ & \textbf{TEXAS} & \textbf{TEXAS} & \textbf{BAYSPAR} \\
 & & & & (\nitrate=10) & (\nitrate=0.1) & \\
\midrule
\multirow{3}{=}{South Dover Bridge}
 & post-PETM & 3  & 30 & 30.2 (+0.02) & 30.2 (+0.02) & 32.5 (+0.25) \\
 & PETM body & 18 & 35 & 35.5 (+0.05) & 32.9 (\textminus0.21) & 39.7 (+0.47) \\
 & pre-PETM  & 14 & 30 & 28.0 (\textminus0.20) & 28.1 (\textminus0.19) & 27.7 (\textminus0.23) \\
\midrule
\multirow{3}{=}{ODP~959}
 & post-PETM & 4  & 30 & 36.5 (+0.65) & 36.4 (+0.64) & 38.1 (+0.81) \\
 & PETM body & 20 & 35 & 42.3 (+0.73) & 40.3 (+0.53) & 43.1 (+0.81) \\
 & pre-PETM  & 16 & 30 & 34.8 (+0.48) & 34.8 (+0.48) & 35.9 (+0.59) \\
\midrule
ODP~1259 & upper asymptote & 59 & 40 & 40.6 (+0.06) & --- & 41.1 (+0.11) \\
Co1010   & lower asymptote & 36 &  5 &  3.0 (\textminus0.20) & --- &  5.8 (+0.08) \\
\bottomrule
\end{tabular*}

\vspace{8pt}
% ---- (B) width -------------------------------------------------------------
\textbf{(B) One-sided posterior width (\degree C) and prior variance removed}\\[2pt]

\begin{tabular*}{\linewidth}{@{}>{\raggedright\arraybackslash}p{\colRecord} >{\raggedright\arraybackslash}p{\colInterval} >{\centering\arraybackslash}p{\colN} @{\extracolsep{\fill}} l c c c c@{}}
\toprule
\textbf{Record} & \textbf{Interval} & $n$ & \textbf{Method}
 & $\sigma_-$ & $\sigma_+$ & $\bar{\sigma}$ & \textbf{resolved} \\
\midrule
\multirow{9}{=}{South Dover Bridge}
 & \multirow{3}{*}{post-PETM} & \multirow{3}{*}{3}
   & TEXAS (\nitrate=10)  & 1.95 & 1.93 & 1.94 & 96\% \\
 & & & TEXAS (\nitrate=0.1) & 1.94 & 1.89 & 1.91 & 96\% \\
 & & & BAYSPAR              & 3.95 & 4.60 & 4.28 & 82\% \\
\cmidrule(l){2-8}
 & \multirow{3}{*}{PETM body} & \multirow{3}{*}{18}
   & TEXAS (\nitrate=10)  & 2.17 & 2.80 & 2.49 & 94\% \\
 & & & TEXAS (\nitrate=0.1) & 2.20 & 3.05 & 2.62 & 93\% \\
 & & & BAYSPAR              & 4.61 & 5.51 & 5.06 & 74\% \\
\cmidrule(l){2-8}
 & \multirow{3}{*}{pre-PETM} & \multirow{3}{*}{14}
   & TEXAS (\nitrate=10)  & 2.09 & 1.93 & 2.01 & 96\% \\
 & & & TEXAS (\nitrate=0.1) & 2.03 & 1.95 & 1.99 & 96\% \\
 & & & BAYSPAR              & 3.78 & 3.98 & 3.88 & 85\% \\
\midrule
\multirow{9}{=}{ODP~959}
 & \multirow{3}{*}{post-PETM} & \multirow{3}{*}{4}
   & TEXAS (\nitrate=10)  & 2.02 & 2.49 & 2.25 & 95\% \\
 & & & TEXAS (\nitrate=0.1) & 2.07 & 2.57 & 2.32 & 95\% \\
 & & & BAYSPAR              & 4.11 & 4.24 & 4.18 & 83\% \\
\cmidrule(l){2-8}
 & \multirow{3}{*}{PETM body} & \multirow{3}{*}{20}
   & TEXAS (\nitrate=10)  & 3.71 & 5.99 & 4.85 & 76\% \\
 & & & TEXAS (\nitrate=0.1) & 4.09 & 6.73 & 5.41 & 71\% \\
 & & & BAYSPAR              & 4.70 & 4.98 & 4.84 & 77\% \\
\cmidrule(l){2-8}
 & \multirow{3}{*}{pre-PETM} & \multirow{3}{*}{16}
   & TEXAS (\nitrate=10)  & 1.94 & 2.22 & 2.08 & 96\% \\
 & & & TEXAS (\nitrate=0.1) & 1.96 & 2.21 & 2.09 & 96\% \\
 & & & BAYSPAR              & 4.01 & 4.12 & 4.07 & 83\% \\
\midrule
\multirow{2}{=}{ODP~1259} & \multirow{2}{*}{upper asymptote} & \multirow{2}{*}{59}
   & TEXAS (\nitrate=10)  & 3.75 & 7.00 & 5.37 & 71\% \\
 & & & BAYSPAR              & 4.98 & 5.14 & 5.06 & 74\% \\
\midrule
\multirow{2}{=}{Co1010} & \multirow{2}{*}{lower asymptote} & \multirow{2}{*}{36}
   & TEXAS (\nitrate=10)  & 7.33 & 5.85 & 6.59 & 57\% \\
 & & & BAYSPAR              & 5.05 & 4.67 & 4.86 & 76\% \\
\bottomrule
\end{tabular*}

\vspace{2pt}
\begin{minipage}{0.95\linewidth}\footnotesize
Note: $\sigma_-$ and $\sigma_+$ are the one-sided posterior widths below and above the median, $(p_{50}-p_{16})/z$ and $(p_{84}-p_{50})/z$ with $z=\Phi^{-1}(0.84)$; $\bar{\sigma}$ is their mean, the symmetric-equivalent width. \textbf{\emph{Resolved}} is $1-\bar{\sigma}^{2}/\sigma_{\text{prior}}^{2}$, the fraction of the prior variance eliminated by the GDGT data: 100\% means the proxy, not the prior, sets the estimate; 0\% means the posterior is as wide as the prior and the data added nothing. All methods use the same per-sample prior mean $\mu_{\text{T,prior}}$ and $\sigma_{\text{prior}}=10$\,\degree C, and within each interval are averaged over the same samples. \nitrate=0.1 is a sensitivity scenario applied only inside the PETM interval, so it is not run for the two asymptote records.
\end{minipage}
\end{table}

% =====================================================================

% Fallbacks so the table compiles even without textcomp.
\providecommand{\textminus}{-}
% Nitrate scenario label, in text mode so it follows the table font.
\providecommand{\nitrate}{NO\textsubscript{3}\textsuperscript{\textminus}}
% Shared geometry of the first three columns -- panels A and B both use these.
\providecommand{\colRecord}{2.55cm}
\providecommand{\colInterval}{2.20cm}
\providecommand{\colN}{0.45cm}
% NB: the column tokens themselves must be written out in each panel --
% \array does not expand a macro in the preamble-token position. Only the
% p{} widths above are shared, and they are the only thing to edit.

\begin{table}[h]\centering\footnotesize
\fontfamily{qhvc}\selectfont
\caption{Prior-referenced comparison of TEXAS and BAYSPAR.}

% ---- (A) temperature: did the data move the posterior off the prior? --------
\textbf{(A) Posterior median temperature (\degree C), with the shift from the prior in units of $\sigma_{\text{prior}}$}\\[2pt]

\begin{tabular*}{\linewidth}{@{}>{\raggedright\arraybackslash}p{\colRecord} >{\raggedright\arraybackslash}p{\colInterval} >{\centering\arraybackslash}p{\colN} @{\extracolsep{\fill}} c c c c@{}}
\toprule
\textbf{Record} & \textbf{Interval} & $n$ & $\mu_{\text{T,prior}}$ & \textbf{TEXAS} & \textbf{TEXAS} & \textbf{BAYSPAR} \\
 & & & & (\nitrate=10) & (\nitrate=0.1) & \\
\midrule
\multirow{3}{=}{South Dover Bridge}
 & post-PETM & 3  & 30 & 30.2 (+0.02) & 30.2 (+0.02) & 32.5 (+0.25) \\
 & PETM body & 18 & 35 & 35.5 (+0.05) & 32.9 (\textminus0.21) & 39.7 (+0.47) \\
 & pre-PETM  & 14 & 30 & 28.0 (\textminus0.20) & 28.1 (\textminus0.19) & 27.7 (\textminus0.23) \\
\midrule
\multirow{3}{=}{ODP~959}
 & post-PETM & 4  & 30 & 36.5 (+0.65) & 36.4 (+0.64) & 38.1 (+0.81) \\
 & PETM body & 20 & 35 & 42.3 (+0.73) & 40.3 (+0.53) & 43.1 (+0.81) \\
 & pre-PETM  & 16 & 30 & 34.8 (+0.48) & 34.8 (+0.48) & 35.9 (+0.59) \\
\midrule
ODP~1259 & upper asymptote & 59 & 40 & 40.6 (+0.06) & --- & 41.1 (+0.11) \\
Co1010   & lower asymptote & 36 &  5 &  3.0 (\textminus0.20) & --- &  5.8 (+0.08) \\
\bottomrule
\end{tabular*}

\vspace{8pt}
% ---- (B) width -------------------------------------------------------------
\textbf{(B) One-sided posterior width (\degree C) and prior variance removed}\\[2pt]

\begin{tabular*}{\linewidth}{@{}>{\raggedright\arraybackslash}p{\colRecord} >{\raggedright\arraybackslash}p{\colInterval} >{\centering\arraybackslash}p{\colN} @{\extracolsep{\fill}} l c c c c@{}}
\toprule
\textbf{Record} & \textbf{Interval} & $n$ & \textbf{Method}
 & $\sigma_-$ & $\sigma_+$ & $\bar{\sigma}$ & \textbf{resolved} \\
\midrule
\multirow{9}{=}{South Dover Bridge}
 & \multirow{3}{*}{post-PETM} & \multirow{3}{*}{3}
   & TEXAS (\nitrate=10)  & 1.95 & 1.93 & 1.94 & 96\% \\
 & & & TEXAS (\nitrate=0.1) & 1.94 & 1.89 & 1.91 & 96\% \\
 & & & BAYSPAR              & 3.95 & 4.60 & 4.28 & 82\% \\
\cmidrule(l){2-8}
 & \multirow{3}{*}{PETM body} & \multirow{3}{*}{18}
   & TEXAS (\nitrate=10)  & 2.17 & 2.80 & 2.49 & 94\% \\
 & & & TEXAS (\nitrate=0.1) & 2.20 & 3.05 & 2.62 & 93\% \\
 & & & BAYSPAR              & 4.61 & 5.51 & 5.06 & 74\% \\
\cmidrule(l){2-8}
 & \multirow{3}{*}{pre-PETM} & \multirow{3}{*}{14}
   & TEXAS (\nitrate=10)  & 2.09 & 1.93 & 2.01 & 96\% \\
 & & & TEXAS (\nitrate=0.1) & 2.03 & 1.95 & 1.99 & 96\% \\
 & & & BAYSPAR              & 3.78 & 3.98 & 3.88 & 85\% \\
\midrule
\multirow{9}{=}{ODP~959}
 & \multirow{3}{*}{post-PETM} & \multirow{3}{*}{4}
   & TEXAS (\nitrate=10)  & 2.02 & 2.49 & 2.25 & 95\% \\
 & & & TEXAS (\nitrate=0.1) & 2.07 & 2.57 & 2.32 & 95\% \\
 & & & BAYSPAR              & 4.11 & 4.24 & 4.18 & 83\% \\
\cmidrule(l){2-8}
 & \multirow{3}{*}{PETM body} & \multirow{3}{*}{20}
   & TEXAS (\nitrate=10)  & 3.71 & 5.99 & 4.85 & 76\% \\
 & & & TEXAS (\nitrate=0.1) & 4.09 & 6.73 & 5.41 & 71\% \\
 & & & BAYSPAR              & 4.70 & 4.98 & 4.84 & 77\% \\
\cmidrule(l){2-8}
 & \multirow{3}{*}{pre-PETM} & \multirow{3}{*}{16}
   & TEXAS (\nitrate=10)  & 1.94 & 2.22 & 2.08 & 96\% \\
 & & & TEXAS (\nitrate=0.1) & 1.96 & 2.21 & 2.09 & 96\% \\
 & & & BAYSPAR              & 4.01 & 4.12 & 4.07 & 83\% \\
\midrule
\multirow{2}{=}{ODP~1259} & \multirow{2}{*}{upper asymptote} & \multirow{2}{*}{59}
   & TEXAS (\nitrate=10)  & 3.75 & 7.00 & 5.37 & 71\% \\
 & & & BAYSPAR              & 4.98 & 5.14 & 5.06 & 74\% \\
\midrule
\multirow{2}{=}{Co1010} & \multirow{2}{*}{lower asymptote} & \multirow{2}{*}{36}
   & TEXAS (\nitrate=10)  & 7.33 & 5.85 & 6.59 & 57\% \\
 & & & BAYSPAR              & 5.05 & 4.67 & 4.86 & 76\% \\
\bottomrule
\end{tabular*}

\vspace{2pt}
\begin{minipage}{0.95\linewidth}\footnotesize
Note: $\sigma_-$ and $\sigma_+$ are the one-sided posterior widths below and above the median, $(p_{50}-p_{16})/z$ and $(p_{84}-p_{50})/z$ with $z=\Phi^{-1}(0.84)$; $\bar{\sigma}$ is their mean, the symmetric-equivalent width. \textbf{\emph{Resolved}} is $1-\bar{\sigma}^{2}/\sigma_{\text{prior}}^{2}$, the fraction of the prior variance eliminated by the GDGT data: 100\% means the proxy, not the prior, sets the estimate; 0\% means the posterior is as wide as the prior and the data added nothing. All methods use the same per-sample prior mean $\mu_{\text{T,prior}}$ and $\sigma_{\text{prior}}=10$\,\degree C, and within each interval are averaged over the same samples. \nitrate=0.1 is a sensitivity scenario applied only inside the PETM interval, so it is not run for the two asymptote records.
\end{minipage}
\end{table}
